% page 617 du fac-simile (image p-616 du scan)
% bloc 172.7 x 246.3 mm ; marge gauche 29.3 mm
\pgc{7.620mm}{0.000mm}{
\l{}
\l{\cel{57}609}
\l{}
\l{\sou{vango}. I. (\sou{anat.}) Che la homo : parto di la vizajo, singla-latere,}
\l{\cel{3}de-sub la okulo til la mentono. - II. Parto korespondanta di}
\l{\cel{3}la kapo, inter la nazo, la boko e la oreli (che la mamiferi) ;}
\l{\cel{3}inter la bazo di la beko e la boko (che la uceli e la fishi) ;}
\l{\cel{3}inter la okulo e la mandibulo (che la insekti). - III. (\sou{tekn}.)}
\l{\cel{3}Parto di objekto plu o min neprecise komparebla, per lua formo}
\l{\cel{3}e simetreso, a la vangi. - D.}
\l{}
\l{\sou{vanilo}. (\sou{bot.)}\cel{3}Frukto di planto sarmentoza di Mexikia. L.\cel{1}\sou{vanilla},}
\l{\cel{3}di qua la shelo odoras aromatoze. - DEFIRS.}
\l{}
\l{\sou{vanitato}. La deziro ostentar. EFIS.}
\l{}
\l{\sou{vanto}. (\sou{navig.})\cel{2}Grosa kordego skalatra qua fixigas masto e mantenas}
\l{\cel{3}lu vertikale. - DR.}
\l{}
\l{\sou{vaporo}. Asemblajo ek la guteti, preske neperceptebla, qui su elevas}
\l{\cel{3}de la surfaco di la liquidi, de la korpi humida, per la efiko}
\l{\cel{3}da kaloro, da la diminuto di la preso atmosferala. - EFIS.}
\l{}
\l{\sou{varo}. To quo esas la objekto di komerco. - DE.}
\l{}
\l{\sou{varango}. (\sou{navig.)}\cel{3}Ligno-peco kurvatra, fixigita per lua mezo,}
\l{\cel{3}sur la kilio di navo, e qua formacas la fundamento di la span-}
\l{\cel{3}taro ye qua konsistas la karpenturo di la navo. - DFS .}
\l{}
\l{\sou{variar}. (\sou{netrans.})\cel{2}Submisar (serio de fakti) a chanji intersuced-}
\l{\cel{3}anta. - II. (\sou{netrans.})\cel{2}Chanjar intersucedante. - Dicesas pri}
\l{\cel{3}kozi qui, kun traiti komuna, interdiferas. - DEFIS.}
\l{}
\l{\sou{varicelo}. (\sou{patol.})\cel{2}Morbo infektala, kontagiala, qua aparas preci-}
\l{\cel{3}pue che la infanti e la pueri, e qua karakterizesas esencale}
\l{\cel{3}da erupto bulo-forma, intervaloze, e qua desaparas pos kelka}
\l{\cel{3}dii. - DEFIS.}
\l{}
\l{\sou{varietato}.\cel{2}(\sou{zool. e bot.}) Singla ek la grupi interdiferanta de}
\l{\cel{3}individui ek un speco, qui havas ula traiti extera komuna. -}
\l{\cel{3}DEFIS.}
\l{}
\l{\sou{variko}.\cel{2}(\sou{patol}.)Dilaturo permananta, quan produktas akumulajo}
\l{\cel{3}de sango en veino (en la gambi, maxim-multa-kaze). - EFIS.}
\l{}
\l{\sou{variolo}. (\sou{patol.})\cel{2}Morbo infektala, kontagiala ed epidemiala, quan}
\l{\cel{3}genitas mikroparazito ankore nekonocata, e qua havas kom}
\l{\cel{3}karakterizivo erupto postuletoza qua abutas a pusifo. - EFIS.}
\l{}
\l{\sou{varma}. I. Di qua la temperaturo esas alta. - II. Qua igas parto}
\l{\cel{3}di la kapo sentar su kaloroza. - DE.}
\l{}
\l{\sou{varsar}. (trans.)\cel{2}Igar (liquido, sablo, grani) fluar, glitar, per}
\l{\cel{3}inklinar la vazo qua kontenas li. - II. (\sou{financo)}. Igar pekunio}
\l{\cel{3}pasar, page, de kaso aden altra kaso. - FI.}
\l{}
\l{\sou{vartar}. (\sou{trans.}) Permanar til la arivo (di ulu, di ulo).- DE.}
}
